\documentclass[11pt,a4paper]{article}
\usepackage[margin=23mm]{geometry}
\usepackage[T1]{fontenc}
\usepackage{lmodern,microtype,amsmath,amssymb,amsthm,mathtools,enumitem,xcolor,booktabs}
\usepackage[colorlinks=true,linkcolor=blue!45!black,citecolor=blue!45!black,urlcolor=blue!45!black]{hyperref}
\newtheorem{theorem}{Theorem}[section]
\newtheorem{lemma}[theorem]{Lemma}
\newtheorem{proposition}[theorem]{Proposition}
\newtheorem{corollary}[theorem]{Corollary}
\theoremstyle{remark}\newtheorem{remark}[theorem]{Remark}
\setlength{\parskip}{3pt}
\setlist{nosep}
\allowdisplaybreaks[1]
\title{Mass gaps, radiation costs and coherent motion\\for mass-critical dispersive equations}
\author{Research continuation: analytic proofs and internal independent checks}
\date{7 October 2026}
\begin{document}
\maketitle
\begin{abstract}
We extend the negative-energy traveling-profile mass gap from biharmonic dispersion to every smooth, real, even, positive elliptic homogeneous symbol of degree greater than one. The proof includes a nonlocal cutoff commutator and Pohozaev identity at energy regularity, and the gap persists under sufficiently small uniform perturbations of the angular symbol. We determine the exact free-radiation energy cost at prescribed mass and momentum: the correct coefficient is the convex envelope of the symbol, even when the dispersion is nonconvex. These results give a mass exclusion for finite traveling-core and dispersive-remainder asymptotics. For the isotropic fourth-order equation, including a nonnegative second-order term, we additionally prove a one-sided center-drift restriction under mass tightness and a first time-average kinetic bound. Failure of the required drift forces divergence of the full initial-time kinetic averages. No soliton-resolution theorem or unrestricted dynamical mass gap is asserted.
\end{abstract}
\noindent\textbf{Status and independence.} This manuscript records complete analytic arguments and internal independent reviews. Its conclusions do not use numerical profiles or the computational certificates of earlier installments. Internal checks are not external peer review; publication priority has not been established. The mass gap proved here is qualitative, with no numerical value of its constant.

\section*{What is being generalized}
The central equation is
\[
 i\partial_t u=a(D)u-|u|^{2m/N}u,\qquad m>1,
\]
where $a$ is real, even, homogeneous of degree $m$, smooth away from the origin, and comparable to $|\xi|^m$. The usual mass-critical BNLS corresponds to $a(\xi)=|\xi|^4$. Throughout, $M(f)=\|f\|_2^2$ denotes mass, not its square root, and the Fourier transform is unitary. A traveling profile has the convention $u(t,x)=e^{i\omega t}U(x-vt)$.

There are three distinct conclusions. The profile gap is an unconditional theorem about every weak profile in the stated class. The radiation optimization is an unconditional variational theorem, while its application to a nonlinear solution requires the stated asymptotic decomposition. The center-drift result is a dynamical theorem for a global strong solution with uniformly concentrated mass; it does not assume a traveling-wave shape. These distinctions are part of the statements.

\section{A mass gap for homogeneous dispersion}\label{sec:r8-homogeneous}

Fix an integer $N\ge1$, a real order $m>1$, and a real, even symbol
$a\in C^{N+2}(\mathbb R^N\setminus\{0\})$ such that
\begin{equation}\label{eq:r8-symbol}
 a(t\xi)=t^m a(\xi)\quad(t>0),\qquad
 c|\xi|^m\le a(\xi)\le C|\xi|^m
\end{equation}
for positive constants $c,C$. Set
\[
 s=\frac m2,\qquad r=\frac mN,\qquad
 \alpha=2r,\qquad p=2(r+1),
\]
and define
\[
 A_a(f)=\int_{\mathbb R^N}a(\xi)|\widehat f(\xi)|^2\,d\xi,\qquad
 M(f)=\|f\|_2^2,\qquad
 E_a(f)=\frac12 A_a(f)-\frac1p\|f\|_p^p.
\]
The Fourier transform is unitary. The critical mass $M_{c,a}>0$ is defined
by the sharp Gagliardo--Nirenberg inequality
\begin{equation}\label{eq:r8-gn}
 \|f\|_p^p\le\frac p2
 \left(\frac{M(f)}{M_{c,a}}\right)^r A_a(f),
 \qquad f\in H^{m/2}(\mathbb R^N;\mathbb C).
\end{equation}
The kinetic form is equivalent to the homogeneous $H^{m/2}$ norm.
In particular, Sobolev interpolation proves that the sharp constant in
\eqref{eq:r8-gn} is finite and positive.

\begin{theorem}[Mass gap for homogeneous dispersion]\label{thm:r8-homogeneous-gap}
There exists $\varepsilon_{N,m,a}>0$ such that every nonzero weak
$H^{m/2}$ solution of
\begin{equation}\label{eq:r8-profile}
 a(D)U+\omega U+i v\cdot\nabla U-|U|^{2m/N}U=0,
 \qquad (\omega,v)\in\mathbb R\times\mathbb R^N,
\end{equation}
with $E_a(U)<0$ satisfies
\[
 M(U)\ge M_{c,a}+\varepsilon_{N,m,a}.
\]
The gap is uniform in frequency and velocity. No radial symmetry,
finite spatial moment, optimizer uniqueness, or real-part
nondegeneracy assumption is required.
\end{theorem}

Thus the exclusion statement applies to all fractional Laplacians
$a(D)=(-\Delta)^{m/2}$ of order $m>1$, all higher integer orders, and
all real-even homogeneous anisotropies satisfying \eqref{eq:r8-symbol}.
It does not assert existence of negative-energy traveling waves for
every such operator. The sharp mass and the gap depend on the symbol.

\subsection{Regularity at the energy level}

\begin{lemma}[Profile regularity]\label{lem:r8-regularity}
Every weak $H^{m/2}$ profile in \eqref{eq:r8-profile} belongs to
$L^\infty\cap H^{m+1}$.
\end{lemma}
\begin{proof}
Since $m>1$, for the fixed velocity $v$ one can choose $\mu>0$ so that
\[
 a(\xi)-v\cdot\xi+\mu\asymp1+|\xi|^m.
\]
The resolvent
\[
 R_\mu=(a(D)+iv\cdot\nabla+\mu)^{-1}
\]
maps $L^q$ to the Bessel potential space $H^{m,q}$ for every
$1<q<\infty$. Indeed, the multiplier
\[
 \frac{(1+|\xi|^2)^{m/2}}
 {a(\xi)-v\cdot\xi+\mu}
\]
satisfies the Mihlin derivative bounds. The possible nonsmoothness
of $a$ at zero is harmless: homogeneity gives the required bounds
after multiplication of each derivative by the corresponding power
of $|\xi|$. The stated $C^{N+2}$ regularity away from zero suffices.
None of the regularity constants below must be uniform in $v$.

Write
\begin{equation}\label{eq:r8-resolvent}
 U=R_\mu(|U|^\alpha U)+(\mu-\omega)R_\mu U.
\end{equation}
We first obtain $U\in L^\infty$. If $N\le m$, the energy space
embeds into all finite $L^q$; it embeds directly into $L^\infty$
when $N<m$. In the endpoint case take a sufficiently large finite $q$.
If $N>m$, start with
\[
 q_0=\frac{2N}{N-m},\qquad
 \beta_0=\frac1{q_0}=\frac12-\frac{m}{2N}.
\]
For $U\in L^{q_j}$ the nonlinear term is in
$L^{q_j/(1+\alpha)}$. Sobolev embedding in
\eqref{eq:r8-resolvent} improves reciprocal integrability according to
\[
 \beta_{j+1}=(1+2m/N)\beta_j-m/N
\]
as long as the right side is positive. The linear resolvent term has
the better reciprocal output exponent $\beta_j-m/N$ and also remains
in $L^2$; interpolation therefore gives the desired common output
exponent. The ideal iterates are
\[
 \beta_j=\frac12-\frac{m}{2N}(1+2m/N)^j.
\]
All nonlinear input exponents exceed one, since
$q_0>1+2m/N$ and the $q_j$ increase. Intermediate exponents can be
weakened slightly to avoid an endpoint. After finitely many steps,
$q_j/(1+\alpha)>N/m$, so the resolvent and Sobolev embedding give
$U\in L^\infty$.

Now $|U|^\alpha U\in L^2$, and \eqref{eq:r8-resolvent} gives
$U\in H^m\subset H^1$. The map
$F(z)=|z|^\alpha z$ on $\mathbb C\simeq\mathbb R^2$ is $C^1$
with bounded derivative on bounded sets. The Sobolev chain rule
therefore gives $F(U)\in H^1$. A final application of the resolvent
yields $U\in H^{m+1}$.
\end{proof}

\subsection{A nonlocal dilation identity without spatial moments}

Let $\chi\in C_c^\infty(\mathbb R^N;\mathbb R)$ equal one near zero, and
put
\[
 f_R(x)=x\chi(x/R)=R f(x/R),\qquad
 L_R=f_R\cdot\nabla+\frac12\operatorname{div}f_R,\qquad R\ge1.
\]
The operator $L_R$ is skew-symmetric on $L^2$.

\begin{lemma}[Cutoff commutator]\label{lem:r8-commutator}
The commutators $[a(D),L_R]$ are uniformly bounded from
$H^{m/2}$ to $H^{-m/2}$, and for every $u\in H^{m/2}$,
\begin{equation}\label{eq:r8-commutator-limit}
 \langle u,[a(D),L_R]u\rangle\longrightarrow mA_a(u)
 \qquad(R\to\infty).
\end{equation}
\end{lemma}
\begin{proof}
Apart from the fixed Fourier normalization factor, the kernel is
\[
 \frac i2\bigl(a(\xi)-a(\eta)\bigr)
 (\xi+\eta)\cdot\widehat f_R(\xi-\eta).
\]
Homogeneity and $C^1$ regularity on the sphere give
\[
 |a(\xi)-a(\eta)|
 \le C|\xi-\eta|(|\xi|+|\eta|)^{m-1}.
\]
For $k=\xi-\eta$, it follows that
\[
 \frac{|a(\xi)-a(\eta)|\,|\xi+\eta|}
 {\langle\xi\rangle^{m/2}\langle\eta\rangle^{m/2}}
 \le C|k|\langle k\rangle^{m/2}.
\]
Here $(|\xi|+|\eta|)^m$ is bounded by a constant times
$\langle\xi\rangle^m+\langle\eta\rangle^m$, and
$\langle\xi\rangle/\langle\eta\rangle$ and its reciprocal are bounded
by a constant times $\langle k\rangle$.
Since $\widehat f_R(k)=R^{N+1}\widehat f(Rk)$,
\[
 \int |k|\langle k\rangle^{m/2}|\widehat f_R(k)|\,dk
 =
 \int |\ell|\langle\ell/R\rangle^{m/2}|\widehat f(\ell)|\,d\ell
 \le C_f.
\]
Young's inequality proves the asserted uniform form bound.

For Schwartz $u$, $L_Ru\to Lu$ in $L^2$, where
$L=x\cdot\nabla+N/2$. Therefore
\[
 \langle u,[a(D),L_R]u\rangle
 =2\operatorname{Re}\langle a(D)u,L_Ru\rangle
 \longrightarrow
 2\operatorname{Re}\langle a(D)u,Lu\rangle=mA_a(u).
\]
The last equality follows from Fourier integration by parts and
$\xi\cdot\nabla a(\xi)=ma(\xi)$. The uniform form bound and
Schwartz density extend this limit to every $u\in H^{m/2}$.
\end{proof}

The same kernel argument, applied to the degree-one symbol
$-v\cdot\xi$, gives
\begin{equation}\label{eq:r8-drift-limit}
 \langle u,[iv\cdot\nabla,L_R]u\rangle
 \longrightarrow-v\cdot P(u).
\end{equation}
For this assertion the form estimate uses $H^{1/2}$, which contains
the energy space continuously because $m>1$.

\begin{lemma}[Profile identities]\label{lem:r8-identities}
Every profile in \eqref{eq:r8-profile} satisfies
\begin{align}
 v\cdot P(U)&=2mE_a(U),\label{eq:r8-pohozaev}\\
 \omega M(U)&=rA_a(U)+(2m-p)E_a(U).\label{eq:r8-frequency}
\end{align}
Writing $U=A+iB$ with real $A,B$, one also has
\begin{equation}\label{eq:r8-phase-energy}
 E_a(U)=\frac1m
 \left\langle B,\bigl(a(D)+\omega-|U|^{2r}\bigr)B\right\rangle.
\end{equation}
\end{lemma}
\begin{proof}
Lemma~\ref{lem:r8-regularity} makes $L_RU\in H^{m/2}$ an admissible
test in the weak equation. The kinetic and drift pairings converge
by \eqref{eq:r8-commutator-limit} and \eqref{eq:r8-drift-limit}.
The frequency pairing vanishes. For the nonlinearity, integration
by parts gives
\[
 \operatorname{Re}\int |U|^\alpha U\,\overline{L_RU}
 =
 \left(\frac12-\frac1p\right)
 \int\operatorname{div}f_R\,|U|^p
 \longrightarrow\frac mp\|U\|_p^p.
\]
This integration is justified by $|U|^\alpha U\in L^2$,
$\nabla U\in L^2$, and the compact support of $f_R$.
The coefficients $\operatorname{div}f_R$ are uniformly bounded
and converge to $N$. The resulting identity is
\[
 \frac m2 A_a(U)-\frac12v\cdot P(U)-\frac mp\|U\|_p^p=0,
\]
which proves \eqref{eq:r8-pohozaev}. Pairing the profile equation
with $U$ gives
$A_a(U)+\omega M(U)-v\cdot P(U)-\|U\|_p^p=0$,
and then \eqref{eq:r8-frequency}.

Because the symbol is real and even, $a(D)$ preserves real
functions. The imaginary profile equation is
\[
 \bigl(a(D)+\omega-|U|^{2r}\bigr)B=-v\cdot\nabla A.
\]
The identity $P(U)=-2\int B\nabla A$ and
\eqref{eq:r8-pohozaev} yield \eqref{eq:r8-phase-energy}.
\end{proof}

\subsection{Compactness and phase coercivity}

\begin{lemma}[Compactness of normalized near-optimizers]
\label{lem:r8-compactness}
The sharp inequality \eqref{eq:r8-gn} has an optimizer.
If $M(f_n)\to M_{c,a}$, $A_a(f_n)=1$, and $E_a(f_n)<0$,
then after translations and passage to a subsequence,
$f_n$ converges strongly in $H^{m/2}$ to a normalized optimizer.
\end{lemma}
\begin{proof}
The sharp inequality implies
$\|f_n\|_p^p\to p/2$ and $E_a(f_n)\to0$.
Local Sobolev nonvanishing gives translations for which a subsequence
has a nonzero weak $H^{m/2}$ limit $Q$. Local Rellich compactness
also gives almost-everywhere convergence after a subsequence.
Write
\[
 m_1=M(Q)>0,\quad a_1=A_a(Q)>0,\quad
 m_2=M_{c,a}-m_1,\quad a_2=1-a_1.
\]
Hilbert norm splitting, the Brezis--Lieb lemma, and
\eqref{eq:r8-gn} give
\[
 M_{c,a}^r\le m_1^ra_1+m_2^ra_2,\qquad a_1+a_2=1.
\]
If $0<m_1<M_{c,a}$, both coefficients $m_j^r$ are strictly below
$M_{c,a}^r$, which is impossible. Hence $m_1=M_{c,a}$, and the
same inequality forces $a_1=1$. Strong energy-space convergence
and equality in the sharp inequality follow.
To prove attainment, normalize an arbitrary maximizing sequence
by amplitude and dilation so that $M(f_n)=M_{c,a}$ and
$A_a(f_n)=1$. Its nonlinear norms converge to $p/2$, and the
identical splitting argument applies; negative energy is not
needed once that convergence is known.
\end{proof}

\begin{lemma}[Constant phase of optimizers]\label{lem:r8-constant-phase}
Every complex optimizer is a constant phase times a real function.
\end{lemma}
\begin{proof}
Write $f=A+iB$, where $A,B$ are real. Evenness of $a$ gives
$A_a(f)=A_a(A)+A_a(B)$. With
$C=(p/2)M_{c,a}^{-r}$, Minkowski, the sharp inequality on each real
component, and two-term H\"older yield
\begin{align*}
 \|f\|_p^2
 &\le\|A\|_p^2+\|B\|_p^2\\
 &\le C^{2/p}\left[
 A_a(A)^{2/p}M(A)^{1-2/p}
 +A_a(B)^{2/p}M(B)^{1-2/p}\right]\\
 &\le C^{2/p}A_a(f)^{2/p}M(f)^{1-2/p}.
\end{align*}
For an optimizer all inequalities are equalities. Strict convexity
of $L^{p/2}$ implies $B^2=cA^2$ almost everywhere, unless one
component vanishes. Repeat this equality argument after multiplying
$f$ by $e^{i\pi/4}$. If both signs of $B/A=\pm\sqrt c$ occurred
on sets of positive measure, the new squared components could
not be proportional. When $c=1$ the new components instead have
disjoint nonzero supports, again giving strict Minkowski inequality.
Thus only one sign is possible, and $f$ has constant phase.
\end{proof}

Let $Q$ be a real optimizer normalized by $M(Q)=M_{c,a}$ and
$A_a(Q)=1$. Its Euler--Lagrange equation is
\[
 a(D)Q+\omega_QQ-|Q|^{2r}Q=0,\qquad
 \omega_Q=r/M_{c,a}>0.
\]

\begin{lemma}[Strict phase coercivity]\label{lem:r8-phase-coercivity}
The real quadratic form
\[
 L_-=a(D)+\omega_Q-|Q|^{2r}
\]
is nonnegative, has kernel $\operatorname{span}_{\mathbb R}\{Q\}$,
and obeys
\[
 \langle b,L_-b\rangle\ge c_Q\|b\|_{H^{m/2}}^2
 \quad\text{if }b\text{ is real and }\langle b,Q\rangle=0,
\]
for some $c_Q>0$.
\end{lemma}
\begin{proof}
Retain $A=A_a(Q)$ and $M=M(Q)$ in the algebra, so that
$\omega_Q=rA/M$. For real $b\in H^{m/2}$ set
$B=M(b)$ and $C_1=A_a(b)$.
The sharp inequality applied to $Q+i\sqrt t\,b$ gives
\begin{equation}\label{eq:r8-kernel-comparison}
 F(t):=\int(Q^2+tb^2)^{r+1}
 \le(r+1)(1+tB/M)^r(A+tC_1),\qquad t\ge0.
\end{equation}
The right derivatives at zero prove $L_-\ge0$.
If $b$ belongs to its form kernel, the first derivatives agree, and
\[
 \int |Q|^{2r}b^2=C_1+rAB/M.
\]
The pointwise second-order remainder for $F$ is nonnegative.
Comparison of second-order remainders in
\eqref{eq:r8-kernel-comparison} and Fatou's lemma yield
\[
 D:=\int |Q|^{2r-2}b^4
 \le(r-1)AB^2/M^2+2BC_1/M.
\]
On the zero set of $Q$, the integrand is interpreted through
the remainder: it is infinite if $r<1$ and $b\ne0$, equals $b^4$
if $r=1$, and equals zero if $r>1$.
Cauchy--Schwarz gives
\[
 D\ge\frac{(C_1+rAB/M)^2}{(r+1)A}.
\]
Combining the inequalities yields
$(C_1-AB/M)^2\le0$. Consequently all inequalities are equalities.
Equality in Cauchy--Schwarz implies
$b^2=(B/M)Q^2$ on $\{Q\ne0\}$. Its integral already equals $B$,
so $b=0$ on the remaining zero set as well.
The form kernel is a linear space containing $Q$. Apply the same
conclusion to $b+Q$ and subtract the identity for $b^2$.
It follows that $b$ is one constant multiple of $Q$.

Finally, multiplication by $|Q|^{2r}$ is compact from $H^{m/2}$
to $H^{-m/2}$. Indeed $|Q|^{2r}\in L^{p/(2r)}$,
H\"older with two $L^p$ factors makes the tail small, and the
embedding into local $L^p$ is compact because $p$ lies strictly
below the critical Sobolev exponent. Thus $L_-$ is a compact
form perturbation of the coercive form $a(D)+\omega_Q$.
Nonnegativity and its one-dimensional kernel give the asserted
coercivity on the $L^2$ orthogonal complement of $Q$.
\end{proof}

\subsection{The mass-gap contradiction}

\begin{proof}[Proof of Theorem~\ref{thm:r8-homogeneous-gap}]
If no uniform gap existed, there would be profiles $U_n$ with
$E_a(U_n)<0$ and $M(U_n)\downarrow M_{c,a}$.
The scaling
\[
 U(x)\mapsto\lambda^{N/2}U(\lambda x),\qquad
 (\omega,v)\mapsto(\lambda^m\omega,\lambda^{m-1}v)
\]
preserves mass and energy sign, so normalize $A_a(U_n)=1$.
Lemmas~\ref{lem:r8-compactness} and \ref{lem:r8-constant-phase}
give, after translations and phases,
$U_n\to Q$ strongly in $H^{m/2}$, with $Q$ a real normalized
optimizer. Identity~\eqref{eq:r8-frequency} gives
$\omega_n\to r/M_{c,a}=\omega_Q$.
A further phase tending to zero ensures
$\langle\operatorname{Im}U_n,Q\rangle=0$.

Strong $L^p$ convergence implies
$|U_n|^{2r}\to|Q|^{2r}$ in $L^{p/(2r)}$.
By H\"older with two $L^p$ factors,
\[
 a(D)+\omega_n-|U_n|^{2r}\longrightarrow L_-
 \quad\text{in form norm }H^{m/2}\to H^{-m/2}.
\]
Write $B_n=\operatorname{Im}U_n$. Phase coercivity and
\eqref{eq:r8-phase-energy} imply, for large $n$,
\[
 E_a(U_n)
 =\frac1m\langle B_n,(a(D)+\omega_n-|U_n|^{2r})B_n\rangle
 \ge\frac{c_Q}{2m}\|B_n\|_{H^{m/2}}^2\ge0,
\]
a contradiction. No uniform bound on $v_n$ is needed: the exact
profile identities eliminate velocity from this final argument.
\end{proof}

\subsection{Local robustness in the uniform topology}

\begin{theorem}[Local uniform robustness of the gap]
\label{thm:r8-symbol-robustness}
Fix a symbol $a_0$ satisfying \eqref{eq:r8-symbol}.
There exist $\delta\in(0,1/2)$ and $\varepsilon>0$ such that, for every
real-even homogeneous symbol $a$ of degree $m$, individually
$C^{N+2}$ away from zero and satisfying
\[
 \sup_{\xi\in S^{N-1}}\left|\frac{a(\xi)}{a_0(\xi)}-1\right|
 <\delta,
\]
every negative-energy traveling profile for $a$ satisfies
\[
 M(U)\ge M_{c,a}+\varepsilon.
\]
No uniform bounds on derivatives of $a$, or compactness of this
symbol neighborhood, are required.
\end{theorem}
\begin{proof}
Suppose the assertion fails. Taking the proposed neighborhood size
and gap both to be $1/n$, choose symbols $a_n$ and
negative-energy profiles $U_n$ such that
\[
 \sup_{S^{N-1}}|a_n/a_0-1|<1/n,\qquad
 0<M(U_n)-M_{c,a_n}<1/n.
\]
Normalize $A_{a_n}(U_n)=1$ by critical scaling.
Homogeneity gives
\[
 (1-1/n)A_{a_0}\le A_{a_n}
 \le(1+1/n)A_{a_0}.
\]
Taking sharp constants in \eqref{eq:r8-gn} shows
$M_{c,a_n}\to M_{c,a_0}$.
The sharp inequality and negative energy also give
\[
 \|U_n\|_p^p\to p/2,\qquad E_{a_n}(U_n)\to0,\qquad
 A_{a_0}(U_n)\to1.
\]
The proof of Lemma~\ref{lem:r8-compactness}, with these converging
normalizations, applies to the limiting symbol. Hence, after
translations and phases, $U_n\to Q$ strongly in $H^{m/2}$,
where $Q$ is a real normalized optimizer for $a_0$.
The frequency identity gives $\omega_n\to r/M_{c,a_0}$.

The kinetic forms converge in norm from $H^{m/2}$ to $H^{-m/2}$
by uniform convergence of the symbols on the sphere.
Together with convergence of the nonlinear potentials and
frequencies, this transfers the strict phase coercivity for $Q$
to the sequence. After arranging
$\langle\operatorname{Im}U_n,Q\rangle=0$, the energy identity
\eqref{eq:r8-phase-energy} yields the same contradiction as before.
Profile regularity and the cutoff identity were established separately
for each symbol and each velocity; their constants never entered this
compactness argument. Uniform derivative bounds are therefore
unnecessary.
\end{proof}

\begin{proposition}[Compact families of symbols]\label{thm:r8-symbol-family}
Fix $N\ge1$ and $m>1$. Let $\mathcal K$ be a family of symbols
as in \eqref{eq:r8-symbol}, each individually real, even, and
$C^{N+2}$ away from zero, whose restrictions to $S^{N-1}$ form
a compact subset of $C^0(S^{N-1})$.
There is a single $\varepsilon_{N,m,\mathcal K}>0$ such that every
negative-energy traveling profile for every $a\in\mathcal K$
satisfies $M(U)\ge M_{c,a}+\varepsilon_{N,m,\mathcal K}$.
\end{proposition}
\begin{proof}
The neighborhoods supplied by
Theorem~\ref{thm:r8-symbol-robustness} form an open cover of
$\mathcal K$ in the uniform topology. Choose a finite subcover
and take the minimum of its positive gap constants.
\end{proof}

\begin{remark}
The mechanism is a restriction on coherent negative-energy motion:
near the sharp mass threshold, variational compactness forces a
profile close to a real optimizer, whose only null phase direction is
the constant phase. Removing that direction makes the phase form
strictly positive, whereas a negative-energy traveling profile requires
a negative value of the same form. This argument is independent of
radiality and does not use a maximum principle.
\end{remark}

\section{The exact radiation cost and asymptotic mass thresholds}

In this section $m>1$ denotes the order of dispersion, while $\mu$ denotes a mass. Let
\[
 a\in C^\infty(\mathbb R^N\setminus\{0\};\mathbb R),\qquad
 a(-\xi)=a(\xi),\qquad a(t\xi)=t^m a(\xi)\quad(t>0),
 \qquad c|\xi|^m\le a(\xi)\le C|\xi|^m,
\]
where $c>0$. Extend $a$ continuously by $a(0)=0$. Put $s=m/2$ and
\[
 p=2+\frac{2m}{N},\qquad
 A_a(f)=\int_{\mathbb R^N}a(\xi)|\widehat f(\xi)|^2\,d\xi,
 \qquad E_a(f)=\frac12A_a(f)-\frac1p\|f\|_p^p.
\]
The Fourier transform is unitary, and
\[
 M(f)=\|f\|_2^2,\qquad P(f)=\int_{\mathbb R^N}\xi|\widehat f(\xi)|^2\,d\xi.
\]
All these quantities are continuous on $H^s$: momentum uses $s>1/2$, and the nonlinear exponent is strictly $H^s$-subcritical. Smoothness of $a$ away from zero will be used only for free propagation. The variational result below requires only continuity and positive elliptic homogeneity.

\subsection{Convexification determines the optimal cost}

Let $b=a^{**}$ be the closed convex envelope of $a$. It is even, convex and homogeneous of degree $m$, and
\[
 c|\xi|^m\le b(\xi)\le a(\xi)\le C|\xi|^m.
\]
For homogeneity, observe that $\xi\mapsto t^{-m}b(t\xi)$ is a convex minorant of $a$, and apply the resulting inequality with both $t$ and $t^{-1}$. Let
\[
 B=\{\xi:a(\xi)\le1\},\qquad K=\operatorname{conv}(B).
\]
This is a compact symmetric convex body with nonempty interior. If $\gamma_K$ is its Minkowski functional, then
\begin{equation}\label{eq:r8-convexified-gauge}
 b(\xi)=\gamma_K(\xi)^m.
\end{equation}
Indeed, $\gamma_K^m$ is convex and at most $a$, hence at most $b$. Conversely, $b\le1$ on $B$, and therefore on $K$; homogeneity gives the reverse inequality.

\begin{theorem}[Exact radiation energy at fixed mass and momentum]\label{thm:r8-radiation-cost}
For every $\mu>0$ and $\pi\in\mathbb R^N$,
\[
 \inf\left\{\frac12 A_a(f):f\in H^{m/2},\ M(f)=\mu,\ P(f)=\pi\right\}
 =\frac{b(\pi)}{2\mu^{m-1}}.
\]
The infimum is not attained. Every strictly larger finite energy is attained by a Schwartz function satisfying the two constraints exactly.
\end{theorem}

\begin{proof}
Set $q=\pi/\mu$ and $d\rho=|\widehat f|^2\,d\xi/\mu$. Jensen's inequality for $b\le a$ gives
\[
 \frac{A_a(f)}\mu=\int a\,d\rho\ge\int b\,d\rho\ge b(q)
 =\mu^{-m}b(\pi).
\]
To prove strictness, choose a supporting affine function $\ell$ of $b$ at $q$. Thus $\ell\le b\le a$ and $\ell(q)=b(q)$. Equality in the bound would force $\rho$ to be supported on $\{a=\ell\}$. On each positive ray $r\theta$, $|\theta|=1$, the nonnegative function
\[
 h_\theta(r)=a(\theta)r^m-r\ell_{\rm lin}(\theta)-\ell(0)
\]
is strictly convex, since
\[
 h_\theta''(r)=m(m-1)a(\theta)r^{m-2}>0\qquad(r>0).
\]
It has at most one zero: two zeros would force a negative value between them. Polar-coordinate Fubini therefore shows that $\{a=\ell\}$ has Lebesgue measure zero. This contradicts the absolute continuity and unit mass of $\rho$. The origin is immaterial; in dimension one the same argument applies to the two half-lines.

For sharpness, first assume $q\ne0$ and set $r_0=\gamma_K(q)$. The point $z=q/r_0$ lies on the boundary of $K$. Choose a supporting hyperplane there. The intersection of that hyperplane with $K$ is the convex hull of its intersection with $B$. Carath\'eodory's theorem in the hyperplane consequently gives
\[
 z=\sum_{j=1}^{L}\theta_j z_j,\qquad
 L\le N,\quad\theta_j>0,\quad\sum_j\theta_j=1,\quad a(z_j)=1.
\]
Here $a(z_j)=1$ because the supporting hyperplane has positive distance from zero: if $a(z_j)<1$, a slightly larger radial multiple of $z_j$ would still belong to $B$ and would violate the supporting inequality. The probability measure with masses $\theta_j$ at $r_0z_j$ has mean $q$ and $a$-moment $r_0^m=b(q)$. At $q=0$, use the unit atom at zero.

Replace each distinct atom by a sufficiently narrow smooth compactly supported Fourier bump. Choose each bump's squared modulus to have unit mass and zero centered first moment, and give it coefficient $\sqrt{\mu\theta_j}$. The supports can be made disjoint. The resulting inverse Fourier transform is Schwartz, has mass $\mu$ and momentum $\pi$ exactly, and its kinetic energy converges to $\mu b(q)/2$. Continuity of $a$ suffices for this convergence.

Finally, fix an energy $e$ above this infimum and choose one such function $f_0$ with energy below $e$. Two equal Fourier bumps centered at $q\pm R\eta$, for a fixed unit vector $\eta$, give another Schwartz function $f_1$ with the same mass and momentum and energy tending to infinity as $R\to\infty$. For sufficiently large $R$ its Fourier support is disjoint from that of $f_0$. The functions
\[
 f_\theta=\sqrt{1-\theta}\,f_0+\sqrt\theta\,f_1,
 \qquad0\le\theta\le1,
\]
have the same mass and momentum, and their kinetic energies interpolate affinely. A suitable $\theta$ gives energy exactly $e$.
\end{proof}

\begin{remark}[An explicit nonconvex elliptic dispersion]
Convexification cannot be omitted for a general elliptic symbol. In two dimensions let
\[
 a(\xi_1,\xi_2)=\xi_1^4+\xi_2^4-\frac32\xi_1^2\xi_2^2
 =\left(\xi_2^2-\frac34\xi_1^2\right)^2+\frac7{16}\xi_1^4.
\]
One has $a(\xi)\ge|\xi|^4/8$. The equal atoms at $(1,\sqrt3/2)$ and $(1,-\sqrt3/2)$ have mean $(1,0)$ and $a$-moment $7/16$. Conversely, every probability measure of that mean has
\[
 \int a\,d\rho\ge\frac7{16}\int\xi_1^4\,d\rho\ge\frac7{16}.
\]
Hence $a^{**}(1,0)=7/16$, whereas $a(1,0)=1$. At momentum $(r,0)$ the sharp cost is $7r^4/(32\mu^3)$, not $r^4/(2\mu^3)$. For the isotropic symbol $a(\xi)=|\xi|^m$, convexity gives $b=a$.
\end{remark}

\subsection{Free propagation and automatic decoupling}

\begin{lemma}[Decay and uniform weak dispersion]\label{lem:r8-free-decay}
Write $S_a(t)=e^{-it a(D)}$. For every $f\in H^{m/2}$ and every finite $q>2$ for which $H^{m/2}$ embeds continuously into $L^q$,
\[
 \|S_a(t)f\|_q\longrightarrow0\qquad(t\to+\infty).
\]
In particular this holds for $q=p=2+2m/N$. For every $f,g\in L^2$,
\[
 \sup_{x\in\mathbb R^N}
 \left|\left\langle S_a(t)f,g(\cdot-x)\right\rangle\right|
 \longrightarrow0.
\]
\end{lemma}

\begin{proof}
First suppose $\widehat f\in C_c^\infty(\mathbb R^N\setminus\{0\})$. At any point $\xi_0$ of its support, put $e=\xi_0/|\xi_0|$. Euler homogeneity gives
\[
 e^T D^2a(\xi_0)e=\frac{m(m-1)}{|\xi_0|^2}a(\xi_0)>0.
\]
This fixed-direction second derivative has a positive lower bound in a small ball about $\xi_0$. Cover the compact Fourier support by finitely many such balls and use a smooth partition of unity. On one piece write $\xi=re+\eta$, with $\eta\perp e$. Along each line intersecting its ball, the Fourier phase obeys
\[
 \Phi(r)=x\cdot(re+\eta)-t a(re+\eta),\qquad
 |\Phi''(r)|\ge c_0|t|.
\]
The second-derivative van der Corput estimate, on the interval cut out by the ball, yields
\[
 \left|\int e^{i\Phi(r)}g(r,\eta)\,dr\right|
 \le C|t|^{-1/2}
 \left(\sup_r|g(r,\eta)|+\int|\partial_rg(r,\eta)|\,dr\right).
\]
Integrating in $\eta$ and summing the finite cover proves
\[
 \|S_a(t)f\|_\infty\le C_f|t|^{-1/2},
\]
uniformly in $x$. No global convexity or nondegeneracy of the full Hessian is required. This estimate has a datum-dependent constant; it is not a uniform $H^{m/2}$-to-$L^\infty$ bound.

Interpolation with the conserved $L^2$ norm gives $L^q$ decay for this class. Functions whose Fourier transforms are smooth, compactly supported, and avoid zero are dense in $H^{m/2}$. Approximation, the isometry of $S_a(t)$ on $H^{m/2}$, and Sobolev embedding give the asserted decay for general data.

For the weak assertion first take $f$ in the same dense class and $g$ Schwartz. Bound the pairing by $\|S_a(t)f\|_\infty\|g\|_1$, uniformly in translation. Approximation of both factors in $L^2$, together with unitarity, completes the argument.
\end{proof}

\begin{theorem}[Decoupling of a strong asymptotic expansion]\label{thm:r8-decoupling}
Suppose a solution with conserved mass, momentum and energy admits
\[
 u(t)=\sum_{j=1}^{J}e^{i\theta_j(t)}U_j(\cdot-x_j(t))
       +S_a(t)f+o_{H^{m/2}}(1),
\]
where $J$ is finite, $U_j,f\in H^{m/2}$, and
$|x_j(t)-x_k(t)|\to\infty$ for $j\ne k$. Then
\begin{align*}
 M(u)&=\sum_jM(U_j)+M(f),\\
 P(u)&=\sum_jP(U_j)+P(f),\\
 E_a(u)&=\sum_jE_a(U_j)+\frac12A_a(f).
\end{align*}
The profiles $U_j$ need not satisfy a traveling-wave equation for this decoupling statement.
\end{theorem}

\begin{proof}
Translation separation eliminates core-core quadratic cross terms by $L^2$ approximation with compactly supported functions. Lemma~\ref{lem:r8-free-decay} eliminates core-radiation cross terms, uniformly in the centers and phases. For the kinetic energy apply its weak assertion to the $L^2$ factors $a(D)^{1/2}f$ and $a(D)^{1/2}U_j$.

For momentum, including the range $1<m<2$, use
\[
 \int\xi_k\overline{\widehat f}\,\widehat g
 =\left\langle |D|^{1/2}f,
 \frac{D_k}{|D|}|D|^{1/2}g\right\rangle.
\]
Both factors lie in $L^2$, and the multiplier $D_k/|D|$ is bounded, with any value assigned at zero. The same separation and weak-dispersion arguments therefore apply.

Approximation of each core by a compactly supported function in $L^p$ decouples their nonlinear norms. The free term tends to zero in $L^p$ by Lemma~\ref{lem:r8-free-decay}. Finally all remainder terms vanish by continuity of the invariants on $H^{m/2}$.
\end{proof}

\subsection{Sharp invariant budgets and the mass gap}

For a decomposition as above, set
\[
 S=\sum_jE_a(U_j),\qquad P_c=\sum_jP(U_j),\qquad\mu=M(f).
\]
If $\mu>0$, total momentum zero and negative total energy imply
\begin{equation}\label{eq:r8-budget}
 \boxed{\quad S<0,\qquad 2\mu^{m-1}(-S)>b(P_c).\quad}
\end{equation}
Indeed, Theorems~\ref{thm:r8-radiation-cost} and~\ref{thm:r8-decoupling} give
\[
 E_a(u)>S+\frac{b(P_c)}{2\mu^{m-1}}.
\]
Conversely, for fixed cores and a mass $\mu$ satisfying~\eqref{eq:r8-budget}, Theorem~\ref{thm:r8-radiation-cost} supplies Schwartz radiation of mass $\mu$, momentum $-P_c$, and free energy smaller than $-S$. If a total energy $E_0$ is prescribed, the exact feasibility condition is
\begin{equation}\label{eq:r8-prescribed-budget}
 E_0-S>\frac{b(P_c)}{2\mu^{m-1}}.
\end{equation}
If $f=0$, the corresponding conditions are simply $P_c=0$ and $S<0$.

\begin{remark}[The extent of sufficiency]
The sufficiency just stated is exact for the radiation datum and the asymptotic invariant budget. It does not construct a nonlinear solution having the prescribed expansion. A nonlinear wave operator or a gluing theorem is a separate requirement. The decoupling theorem also assumes an actual strong expansion and conservation; it is not a well-posedness or soliton-resolution theorem.
\end{remark}

\begin{proposition}[Transfer of the profile gap to asymptotic dynamics]\label{prop:r8-asymptotic-gap}
Assume that every negative-energy traveling profile for the symbol $a$ has mass at least $M_{c,a}+\varepsilon_{N,m,a}$, as supplied by the preceding mass-gap theorem. Let a solution of negative total energy have the expansion in Theorem~\ref{thm:r8-decoupling}, with all $U_j$ traveling profiles. Then
\[
 M(u)\ge M_{c,a}+\varepsilon_{N,m,a}.
\]
If exactly $k$ cores have negative energy, their masses alone contribute at least $k(M_{c,a}+\varepsilon_{N,m,a})$.

If there is no radiation, every core has negative energy, and total momentum is zero, then $J\ge2$ and
\[
 M(u)\ge2(M_{c,a}+\varepsilon_{N,m,a}).
\]
\end{proposition}

\begin{proof}
The radiation energy is nonnegative, so negative total energy requires at least one negative-energy core. Apply the profile mass gap to each such core and use mass decoupling. In the final assertion, a single negative-energy traveling profile cannot have zero momentum, because its Pohozaev identity is $v\cdot P=2m E_a<0$. Thus there must be at least two cores.
\end{proof}

\begin{proposition}[A sequential version with a general dispersive remainder]\label{prop:r8-sequential-gap}
The single-gap conclusion and the count of negative-energy cores remain valid under the weaker assumption that, along a sequence $t_n\to\infty$,
\[
 u(t_n)=\sum_{j=1}^{J}e^{i\gamma_{j,n}}U_j(\cdot-x_{j,n})
          +r_n+o_{H^{m/2}}(1),
\]
where the finitely many centers separate, each $U_j$ is a traveling profile, $r_n$ is bounded in $H^{m/2}$,
$\|r_n\|_p\to0$, and $r_n(\cdot+x_{j,n})\rightharpoonup0$ in $H^{m/2}$ for every $j$. No free-evolution representation of the remainder, full-time resolution, or zero-total-momentum assumption is needed.
\end{proposition}

\begin{proof}
Separation and the assumed weak convergence give quadratic splitting, including for the continuous kinetic form on $H^{m/2}$. Compact-support approximation separates the nonlinear norms of the cores, and $\|r_n\|_p\to0$ eliminates all nonlinear remainder terms. Thus
\[
 M(u)=\sum_jM(U_j)+M(r_n)+o(1),\qquad
 E_a(u)=\sum_jE_a(U_j)+\tfrac12 A_a(r_n)+o(1).
\]
The bounded nonnegative sequence $A_a(r_n)$ has a convergent subsequence. Negative total energy then forces at least one negative-energy core. Apply the profile mass gap to each such core and use the mass identity. If $J=0$, the negative-energy assumption is impossible.
\end{proof}

The doubled lower bound applies to the specified radiation-free class. For a single negative core with radiation, the sharp additional-mass requirement is only
\[
 \mu>\left(\frac{b(P(U))}{2|E_a(U)|}\right)^{1/(m-1)}.
\]
For a hypothetical small-speed branch with
$P(U_v)=v\pi_1+o(v)$, $\pi_1\ne0$, and
$-E_a(U_v)=c_1v^2+o(v^2)$, $c_1>0$, this cost is asymptotic to a positive constant times
\[
 |v|^{(m-2)/(m-1)}.
\]
It tends to zero for $m>2$, remains of constant order for $m=2$, and diverges for $1<m<2$. This is a conditional invariant obstruction; it does not assert existence of negative-energy branches in those orders. Nor does the asymptotic mass gap classify all bounded solutions: extending it beyond the stated decomposition remains a dynamical problem.

\section{Mass localization, arbitrary momentum, and averaged kinetic growth}

We consider
\begin{equation}\label{drift:equation}
 i\partial_tu=\Delta^2u-\mu\Delta u-|u|^{8/N}u,
 \qquad N\ge1,\quad \mu\ge0,
\end{equation}
and standard strong $H^2$ solutions for which mass, momentum, and energy are conserved. Our conventions are
\begin{align*}
 M&=\int |u|^2, & P&=\operatorname{Im}\int\overline u\nabla u,\\
 E_\mu&=\frac12\|\Delta u\|_2^2+\frac\mu2\|\nabla u\|_2^2
       -\frac1{2+8/N}\int|u|^{2+8/N}.
\end{align*}
Write $K(t)=\|\Delta u(t)\|_2$, $h(t)=1+K(t)^2$, and
\[
 W(t)=\operatorname{Im}\int\Delta u\nabla\overline u,
 \qquad V(t)=4W(t)+2\mu P.
\]
Thus $V$ is the total mass flux. It is well defined in $H^2$, even when the spatial first moment is infinite.

The compact-orbit relation between translation, momentum, and energy is already present in Pausader--Shao~\cite[Propositions 6.1--6.2]{PS}; their Remark 6.1 explicitly describes the interaction-virial interpretation. Related $H^2$-compactness arguments appear in~\cite[Proposition 5.2]{Komada}. The extension below requires only tightness of mass, allows noncompact curvature tails, and assumes only a first time average of kinetic energy. Its one-sided conclusion is essential.

\begin{theorem}[Localized mass forces counter-momentum drift]\label{drift:main}
Let $u$ be a forward-global solution of~\eqref{drift:equation}, with $M>0$. Suppose there is a center $c:[0,\infty)\to\mathbb R^N$ such that
\begin{equation}\label{drift:tight}
 m(r):=\sup_{t\ge0}\int_{|x-c(t)|>r}|u(t,x)|^2\,dx\longrightarrow0
 \quad(r\to\infty).
\end{equation}
No regularity of $c$ is assumed. Suppose also that, for some $H,T_0<\infty$ with $T_0>0$,
\begin{equation}\label{drift:budget}
 \sup_{S\ge0,\,T\ge T_0}\frac1T\int_S^{S+T}h(t)\,dt\le H.
\end{equation}
Then
\begin{equation}\label{drift:current}
 \lim_{T\to\infty}\sup_{S\ge0}
 \left|\frac{c(S+T)-c(S)}T
       -\frac1{MT}\int_S^{S+T}V(t)\,dt\right|=0,
\end{equation}
and
\begin{equation}\label{drift:inequality}
 \limsup_{T\to\infty}\sup_{S\ge0}
 \left\{\frac{P\cdot[c(S+T)-c(S)]}T
       +\frac{2\mu}T\int_S^{S+T}\|\nabla u(t)\|_2^2\,dt\right\}
 \le8E_\mu.
\end{equation}
In particular, if $E_\mu<0$, then $P\ne0$ and
\[
 \liminf_{T\to\infty}\inf_{S\ge0}
 \frac{-P\cdot[c(S+T)-c(S)]}{|P|T}
 \ge\frac{8|E_\mu|}{|P|}.
\]
\end{theorem}

\begin{theorem}[Failure of the required drift forces averaged growth]\label{drift:growth}
Let $u$ be a forward-global solution with $E_\mu<0$ satisfying~\eqref{drift:tight}. If, for some $\delta>0$ and all sufficiently large $t$,
\begin{equation}\label{drift:slow}
 P\cdot[c(t)-c(0)]\ge(8E_\mu+\delta)t,
\end{equation}
then
\begin{equation}\label{drift:divergence}
 \lim_{T\to\infty}\frac1T\int_0^T\|\Delta u(t)\|_2^2\,dt=\infty.
\end{equation}
In particular this holds if $c(t)=o(t)$, or if
$\liminf_{t\to\infty}P\cdot c(t)/t\ge0$. No zero-momentum assumption is imposed.
\end{theorem}

\subsection{The interaction estimate with its full mass-flux term}

We give the estimate and its tail closure explicitly. Choose a smooth radial function $\phi_R$ which equals $|x|^2/2$ on $|x|\le4R$, is constant on $|x|\ge8R$, and satisfies
\[
 D^2\phi_R\le I,\qquad
 \|D^j\phi_R\|_\infty\le C_jR^{2-j}\quad(1\le j\le6).
\]
Set $\rho=|u|^2$, $j=\operatorname{Im}(\overline u\nabla u)$, $w=\operatorname{Im}(\Delta u\nabla\overline u)$, and
\begin{align*}
 I_R(t)&=2\iint\rho(t,y)\nabla\phi_R(x-y)\cdot j(t,x)\,dx\,dy,\\
 d_R(t)&=M^{-2}\iint_{|x-y|>R}\rho(t,x)\rho(t,y)\,dx\,dy,
 \qquad \theta=\min(4/N,1/4).
\end{align*}
All constants below may depend on $N,M,\mu$ and the cutoffs, but not on a pointwise kinetic bound.

\begin{lemma}[Mass-tail interaction bound]\label{drift:interaction}
For $R\ge1$,
\begin{align}
 |I_R(t)|&\le C R h(t)^{1/4},\label{drift:endpoint}\\
 I_R'(t)&\le16ME_\mu-4\mu M\|\nabla u\|_2^2
             -8P\cdot W-4\mu|P|^2
             +C R^{-2}h^{1/2}+C h d_R^\theta.
 \label{drift:virial}
\end{align}
The derivative is understood in the integrated energy-space sense.
\end{lemma}

\begin{proof}
Interpolation gives $\|\nabla u\|_2\le G:=M^{1/4}K^{1/2}$, proving~\eqref{drift:endpoint}. At a fixed center $y$, let
$J_R^y=2\int\nabla\phi_R(x-y)\cdot j(x)\,dx$. Direct differentiation and integration by parts give, with repeated indices summed,
\begin{align*}
 (J_R^y)'={}&8\operatorname{Re}\int\partial_{jk}\phi_R\,
                  \partial_{ij}\overline u\,\partial_{ik}u
 -4\operatorname{Re}\int\partial_{jk}\Delta\phi_R\,
                  \partial_j\overline u\,\partial_ku\\
 &-2\int\Delta^2\phi_R|\nabla u|^2+\int\Delta^3\phi_R\rho
 -\frac{2\beta}{\beta+2}\int\Delta\phi_R|u|^{\beta+2}\\
 &+4\mu\operatorname{Re}\int\partial_{jk}\phi_R\,
                    \partial_j\overline u\,\partial_ku
 -\mu\int\Delta^2\phi_R\rho,
 \qquad\beta=8/N.
\end{align*}
Here every cutoff is evaluated at $x-y$. Since $D^2\phi_R\le I$, the leading fourth-order and second-order forms have favorable upper signs. Hence
\begin{equation}\label{drift:single}
 (J_R^y)'\le16E_\mu-4\mu\|\nabla u\|_2^2
       +C R^{-2}h^{1/2}
       +C\int_{|x-y|>4R}|u(x)|^{2+8/N}\,dx.
\end{equation}

The derivative falling on the \emph{second} mass density cannot be discarded. For real $f$ the mass equation is
\begin{equation}\label{drift:mass}
 \frac{d}{dt}\int f\rho
 =2\operatorname{Im}\int(\Delta f)\overline u\Delta u
  +4\int\nabla f\cdot w+2\mu\int\nabla f\cdot j.
\end{equation}
Using $\operatorname{Im}(\overline u\Delta u)=\operatorname{div}j$, its exact contribution to $I_R'$ is
\begin{align}
 Y_R={}&4\iint D^2\Delta\phi_R(x-y)[j(x),j(y)]\,dx\,dy\notag\\
 &-8\iint D^2\phi_R(x-y)[j(x),w(y)]\,dx\,dy\notag\\
 &-4\mu\iint D^2\phi_R(x-y)[j(x),j(y)]\,dx\,dy.
 \label{drift:second-density}
\end{align}
The first term is $O(R^{-2}h^{1/2})$. Replacing the other Hessians by $I$ gives exactly $-8P\cdot W-4\mu|P|^2$.

To bound the errors, write $m_r(a)=\int_{|x-a|>r}\rho(x)\,dx$ and
\[
 H_r=K+C G/r+C\sqrt M/r^2.
\]
For a cutoff $\eta_r$ which vanishes inside radius $r$ and equals one outside radius $2r$, product differentiation, interpolation, and the mass-critical Gagliardo--Nirenberg inequality give
\begin{align}
 \int_{|x-a|>2r}|u|^{2+8/N}&\le C m_r(a)^{4/N}H_r^2,
 \label{drift:nonlinear-tail}\\
 \int_{|x-a|>2r}|\nabla u|^2&\le m_r(a)^{1/2}H_r.
 \label{drift:gradient-tail}
\end{align}
Averaging~\eqref{drift:nonlinear-tail} against $\rho(a)$ and applying Jensen (or $m_r\le M$ when $4/N>1$) bounds the nonlinear remainder in~\eqref{drift:single} by
$C_MH_{2R}^2d_R^{\min(4/N,1)}$.

For the cross-flux error define
\[
 B_R=\iint_{|x-y|>4R}\rho(x)|\nabla u(y)|^2\,dx\,dy.
\]
Equation~\eqref{drift:gradient-tail} and Cauchy--Schwarz imply
\[
 B_R\le H_{2R}\sqrt{M\,M^2d_R}.
\]
Cauchy--Schwarz on the product space therefore gives
\[
 \iint_{|x-y|>4R}|u(x)||\nabla u(x)||\Delta u(y)||\nabla u(y)|
 \le GK\sqrt{B_R}\le C_Mh\,d_R^{1/4}.
\]
The $j$--$j$ tail is at most $B_R$, again by product-space Cauchy--Schwarz and symmetry, and hence at most $C_Mh^{1/2}d_R^{1/2}$. These estimates prove~\eqref{drift:virial}. The calculation initially applies to smooth approximations; its integrated form passes to standard strong $H^2$ solutions because all displayed terms have the indicated continuous energy-space bounds.
\end{proof}

\subsection{Recovering center motion without spatial moments}

Mass tightness implies $d_R(t)\le2m(R/2)/M$. It also determines centers up to a uniformly bounded error: choose $r_0\ge1$ with $m(r_0)\le M/8$, so every $B(c(t),r_0)$ contains at least $7M/8$. This observation also shows that the condition with an infimum over centers at each radius is equivalent to~\eqref{drift:tight}; one first fixes these $7M/8$ centers and compares them with any higher-concentration ball by intersection.

For a fixed cutoff $\chi((x-a)/r)$, $r\ge1$,~\eqref{drift:mass} gives
\begin{equation}\label{drift:cutoff-flux}
 \left|\frac{d}{dt}\int\chi((x-a)/r)\rho\right|
 \le C r^{-1}h(t).
\end{equation}
Under~\eqref{drift:budget}, every interval $J$, including intervals shorter than $T_0$, satisfies
$\int_Jh\le H(|J|+T_0)$. Start a cutoff at $c(s)$, equal to one on radius $L$ and supported on radius $2L$. Taking
$L\ge r_0+C H(t-s+T_0)/M$ retains at least $3M/4$ of its mass until time $t$. Its support intersects $B(c(t),r_0)$, and hence
\begin{equation}\label{drift:coarse}
 |c(t)-c(s)|\le A|t-s|+B
\end{equation}
with uniform constants. More generally, if only $\int_0^Th\le HT$ is known, the same argument gives
\begin{equation}\label{drift:sequence-coarse}
 \sup_{0\le t\le T}|c(t)-c(0)|\le C_HT+C.
\end{equation}

Let $F_R(z)=z\chi(z/R)$ be a smooth vector cutoff, equal to $z$ on $|z|\le R$, zero on $|z|\ge2R$, with bounded first derivatives and second derivatives $O(R^{-1})$. On a window $[S,S+T]$, choose $R=LT$ with $L$ fixed sufficiently large that~\eqref{drift:coarse} puts all centers inside $B(c(S),R/4)$. Define
\[
 Z_R(t)=\int F_R(x-c(S))\rho(t,x)\,dx.
\]
Equation~\eqref{drift:mass} gives $Z_R'=V+\mathcal E_R$. Since the nonconstant-gradient error is supported outside the ball of radius $R$ about $c(S)$,~\eqref{drift:gradient-tail} implies
\begin{equation}\label{drift:moment-error}
 |\mathcal E_R(t)|\le C\left[m(R/4)^{1/4}+m(R/2)^{1/2}\right]h(t)^{3/4}
                         +C R^{-1}h(t)^{1/2}.
\end{equation}
Indeed the fourth-order flux tail is bounded by
$\|\Delta u\|_2(\int_{|x-c(t)|>R/2}|\nabla u|^2)^{1/2}$; the second-order flux tail is bounded by
$m(R/2)^{1/2}\|\nabla u\|_2$. Both have the displayed integrable kinetic dependence.

For either endpoint $t\in\{S,S+T\}$ and any fixed $r$, splitting into the ball about $c(t)$ and its complement yields, for sufficiently large $T$,
\begin{equation}\label{drift:moment-endpoint}
 |Z_R(t)-M[c(t)-c(S)]|\le Mr+C Rm(r).
\end{equation}
The average of~\eqref{drift:moment-error} tends to zero uniformly in $S$, by H\"older and~\eqref{drift:budget}. Divide~\eqref{drift:moment-endpoint} by $T$, first let $T\to\infty$, and then $r\to\infty$. Integration proves~\eqref{drift:current}.

The same argument, using~\eqref{drift:sequence-coarse}, proves the following useful sequence version:
\begin{equation}\label{drift:sequence-current}
 \int_0^{T_n}h\le H T_n,\quad T_n\to\infty
 \quad\Longrightarrow\quad
 c(T_n)-c(0)=\frac1M\int_0^{T_n}V(t)\,dt+o(T_n).
\end{equation}
No center of mass is introduced in either argument.

\subsection{Low-kinetic endpoints and the growth alternative}

\begin{proof}[Proof of Theorem~\ref{drift:main}]
It remains to prove~\eqref{drift:inequality}. Choose
$s\in[S,S+T_0]$ and $t\in[S+T-T_0,S+T]$ with $h(s),h(t)\le2H$, which is possible by~\eqref{drift:budget}. Integrate~\eqref{drift:virial} between these endpoints. Using~\eqref{drift:endpoint}, $d_R\le2m(R/2)/M$, and H\"older gives
\begin{align*}
 2P\cdot\int_s^t V
 +4\mu M\int_s^t\|\nabla u\|_2^2
 \le{}&16ME_\mu(t-s)+C R H^{1/4}\\
 &+C R^{-2}H^{1/2}T
 +C H[2m(R/2)/M]^\theta T.
\end{align*}
The omitted endpoint intervals have uniformly bounded integrals of $|V|$ and $\|\nabla u\|_2^2$, since both are bounded by a constant times $h$ and each interval has length at most $T_0$. Divide by $T$ and let $T\to\infty$, then $R\to\infty$. Finally use~\eqref{drift:current} and divide by $2M$ to obtain~\eqref{drift:inequality}. Its stated negative-energy consequences follow from $\mu\ge0$.
\end{proof}

\begin{proof}[Proof of Theorem~\ref{drift:growth}]
If~\eqref{drift:divergence} fails, there are $T_n\to\infty$ and $H<\infty$ with $\int_0^{T_n}h\le H T_n$. Choose $t_n\in[T_n/2,T_n]$ satisfying $h(t_n)\le4H$. Then
\[
 \int_0^{t_n}h\le2H t_n.
\]
Integrate~\eqref{drift:virial} from $0$ to $t_n$. Both endpoint factors in~\eqref{drift:endpoint} are uniformly bounded for fixed $R$. Divide by $t_n$, pass to the upper limit, and then let $R\to\infty$. Dropping the nonnegative second-order kinetic term gives
\[
 \limsup_n\frac1{M t_n}P\cdot\int_0^{t_n}V(t)\,dt\le8E_\mu.
\]
Equation~\eqref{drift:sequence-current}, applied with budget $2H$, implies
\[
 \limsup_n\frac{P\cdot[c(t_n)-c(0)]}{t_n}\le8E_\mu,
\]
contradicting~\eqref{drift:slow}. This proves divergence of the full averages, not only unboundedness along a subsequence.
\end{proof}

\begin{remark}
The conclusions describe a necessary motion of any sufficiently controlled, mass-localized negative-energy solution; they do not classify such a solution as a traveling wave. They permit arbitrary curvature leakage and, under~\eqref{drift:budget}, arbitrarily large instantaneous kinetic norms. Negative-energy traveling waves satisfy the corresponding exact identity
$P\cdot v+2\mu\|\nabla U\|_2^2=8E_\mu(U)$, while an outgoing pair with zero total momentum fails mass tightness. Thus both known counterexample mechanisms are consistent with these restrictions. For maximal solutions, Theorem~\ref{drift:growth} gives a finite-time-breakdown or averaged-growth alternative under its localization and drift hypotheses; it does not establish finite-time blowup from negative energy alone.
\end{remark}


\section{The remaining unrestricted dynamical question}
For pure BNLS, put $M_c=M_{c,|\xi|^4}$ and consider
\[
 m_{\mathrm{ng}}=\inf\left\{M(u_0):\ E(u_0)<0,\ P(u_0)=0,\quad
 u\text{ is forward global},\quad
 \sup_{t\ge0}\|u(t)\|_{H^2}<\infty\right\}.
\]
Nothing proved here establishes $m_{\mathrm{ng}}>M_c$ without an additional compactness or asymptotic decomposition hypothesis. In particular, the elliptic imaginary-part identity for a traveling profile is not an identity for an arbitrary time-dependent solution.

The obstruction to one tempting argument can be made exact. The sharp inequality supplies only
\[
 0<\frac{|E(f)|}{\|\Delta f\|_2^2}
 \le\frac12\left[\left(\frac{M(f)}{M_c}\right)^{4/N}-1\right]
 \quad\text{when } E(f)<0.
\]
There is no positive lower bound on this ratio from the mass excess alone. Fix any $M_0>M_c$. A real smooth compactly supported function of mass $M_0$ and negative energy is obtained by approximating a suitable multiple of an optimizer and renormalizing. Another real smooth function of the same mass can have positive energy, for example by separating the mass into sufficiently many disjoint subthreshold pieces. Join the two in the mass sphere of a finite-dimensional real smooth compactly supported space. Approaching the first zero of the continuous energy along this path from its negative side gives
\[
 M(f_n)=M_0,\quad P(f_n)=0,\quad E(f_n)<0,\quad
 \frac{|E(f_n)|}{\|\Delta f_n\|_2^2}\longrightarrow0.
\]
The denominator tends to a positive number because the limiting zero-energy function has nonzero mass. These are admissible initial functions; no global behavior for their evolutions is inferred.

Consequently, a small near-optimizer mass tail cannot simply absorb a virial error of the form $C\|\Delta f\|_2^2d_R^\theta$ into $|E|$. To obtain an unrestricted result through this route, a sharper estimate must control the localized interaction cross-flux as well as the nonlinear tail, possibly retaining the favorable Hessian defect discarded in the present upper bound. Such an estimate is not proved here. Neither a general nonlinear wave operator for the prescribed core/radiation data nor general soliton resolution is supplied by conservation-law feasibility.

\section{Relation to existing work}
The constant-phase equality mechanism for complex optimizers has precedent in Proposition 3.8 of~\cite{FJMM}. Related variational phase-kernel and spectral-gap arguments occur in Iyer--Parker--Stefanov~\cite{IPS}, Section 3.3.2, in a one-dimensional fourth-order setting. The proof here spells out the all-dimensional, arbitrary-power algebra used for general homogeneous dispersion; it does not claim to invent phase rotation or the variational method.

Naumkin--Rapha\"el~\cite{NR} construct traveling waves with subthreshold mass for one-dimensional fractional dispersion of order $1<m<2$. There is no conflict: the sharp inequality forces their unmodified energy to be positive. Our gap concerns \emph{negative} energy, and cannot be stated for all traveling waves. General-dispersion boosted-wave existence and symmetry have been studied by Bugiera--Lenzmann--Schikorra--Sok~\cite{BLSS}; our exclusion result does not replace those existence theorems. Fractional radial blowup and nonlocal virial methods also have important prior work~\cite{BHL}.

Pausader--Shao~\cite{PS}, Proposition 6.1, already prove the precise relation between translation, momentum and energy for compact soliton-like solutions. Their Proposition 6.2 and Remark 6.1 identify the averaged-current and interaction-virial connection. Komada~\cite{Komada}, Proposition 5.2, gives a localized interaction identity for bounded $H^2$-precompact trajectories in the mass-supercritical, energy-subcritical setting. The present dynamical increment is the \emph{one-sided} restriction with only mass tightness, a first kinetic time average and arbitrary momentum, together with the full-average divergence consequence. The exact compact-orbit identity is not claimed as new.

Finally, Jensen's inequality, convex relaxation, Carath\'eodory's theorem and van der Corput bounds are standard tools. Their role here is to identify the sharp radiation cost for nonconvex dispersion and establish the decoupling needed for the stated nonlinear corollaries. A search of these closest primary references found no identical general negative-energy profile-gap theorem; that limited comparison is not an exhaustive novelty certification.

\begin{thebibliography}{9}
\bibitem{PS} B. Pausader and S. Shao, \emph{The mass-critical fourth-order Schr\"odinger equation in high dimensions}, \href{https://arxiv.org/abs/0906.1547}{arXiv:0906.1547}, especially Section 6.
\bibitem{Komada} K. Komada, \emph{Scattering of group-invariant solutions below the ground state for the fourth-order NLS}, \href{https://arxiv.org/abs/2609.33381}{arXiv:2609.33381v1} (2026), Proposition 5.2.
\bibitem{NR} I. Naumkin and P. Rapha\"el, \emph{On small travelling waves to the mass critical fractional NLS}, Calc. Var. Partial Differential Equations \textbf{57}, 93 (2018), \href{https://arxiv.org/abs/1712.04484}{arXiv:1712.04484}.
\bibitem{BLSS} L. Bugiera, E. Lenzmann, A. Schikorra and J. Sok, \emph{On symmetry of traveling solitary waves for dispersion generalized NLS}, \href{https://arxiv.org/abs/1908.06771}{arXiv:1908.06771}.
\bibitem{FJMM} A. J. Fern\'andez, L. Jeanjean, R. Mandel and M. Mari\c{s}, \emph{Some non-homogeneous Gagliardo--Nirenberg inequalities and application to a biharmonic non-linear Schr\"odinger equation}, J. Differential Equations \textbf{330} (2022), 1--65, \href{https://arxiv.org/abs/2010.01448}{arXiv:2010.01448}.
\bibitem{IPS} V. Iyer, R. Parker and A. G. Stefanov, \emph{Existence and stability for traveling waves of fourth order semilinear wave and Schr\"odinger equations}, \href{https://arxiv.org/abs/2512.12390}{arXiv:2512.12390}, Section 3.3.2.
\bibitem{BHL} T. Boulenger, D. Himmelsbach and E. Lenzmann, \emph{Blowup for fractional NLS}, \href{https://arxiv.org/abs/1509.08845}{arXiv:1509.08845}.
\end{thebibliography}
\end{document}
